% =============================================================
% Annexe B — Extraits techniques vérifiés dans le dépôt
% =============================================================
\chapter{Extraits techniques de la réalisation}
\label{ann:extraits-realisation}

Cette annexe rassemble les extraits de code que le \cref{ch:realisation} commente sans les
reproduire : courts, localisés, issus du dépôt vérifié le \dateref, et sans aucun secret ni
identifiant de projet réel.

\section{Dépôt d'infrastructure et protection des ressources}

\begin{extraitconf}[boxsep=1pt, left=3pt, right=3pt, top=2pt, bottom=2pt]
terraform/\\
├── environments/\\
│\ \ \ ├── dev/\ \ \ \ \ \ \ \ \ \ \ \ configuration appliquée\\
│\ \ \ ├── staging/\ \ \ \ \ \ \ configuration appliquée, second locataire\\
│\ \ \ └── prod/\ \ \ \ \ \ \ \ \ \ \ cible vide (\texttt{.gitkeep})\\
└── modules/\\
\ \ \ \ ├── app-service, cloud-run, dashboard, workflow\\
\ \ \ \ ├── vpc, load-balancer, cloud-sql, artifact-registry\\
\ \ \ \ ├── iam, kms, audit, logging, monitoring\\
\ \ \ \ └── bigquery, detection, ml-pipeline
\end{extraitconf}

Le découpage transpose les niveaux L1--L7 en modules techniques sans correspondance bijective,
plusieurs contrôles traversant plusieurs modules : \texttt{app-service} crée les ressources propres
à une application, \texttt{iam} porte la fédération, \texttt{bigquery} et \texttt{detection} les
données et les règles, \texttt{ml-pipeline} l'encodage et l'enrichissement. L'état que ce dépôt
produit cartographie le système entier, et c'est à ce titre une donnée sensible : le
\cref{tab:ann-securite-etat} récapitule les cinq protections qui lui sont appliquées.

\begin{table}[!htbp]
\caption{Sécurité du fichier d'état d'infrastructure}
\label{tab:ann-securite-etat}
\centering
\footnotesize
\setlength{\tabcolsep}{4pt}
\begin{tabular}{|P{2.7cm}|P{5.9cm}|P{6.4cm}|}
\hline
\textbf{Élément} & \textbf{Mise en œuvre} & \textbf{Raison} \\
\hline
Emplacement & Stockage objet dédié, versionné et non public & L'état cartographie le système \\ \hline
Séparation & Un préfixe d'état par environnement & Limiter le rayon d'impact d'une erreur \\ \hline
Accès & Identité de la chaîne d'infrastructure et administration & Moindre privilège \\ \hline
Verrouillage & Verrou natif du stockage & Éviter deux applications concurrentes \\ \hline
Secrets & Conteneurs créés en code ; valeurs injectées séparément & Éviter les valeurs sensibles dans l'état \\
\hline
\end{tabular}
\end{table}

Ces cinq mesures réduisent l'exposition de l'état ; leur résistance à une identité non autorisée
est éprouvée par le protocole T17, conforme le 20/08/2026 et rejoué à la cadence du plan de
validation continue (\cref{sec:plan-validation}).


\textbf{Prévention de l'accès public au média.} L'extrait suivant provient de
\texttt{terraform/modules/app-service/main.tf}, lignes 133--143. La protection est imposée sur la
ressource ; elle ne dépend pas d'une politique d'organisation.

\begin{extraitconf}[boxsep=1pt, left=3pt, right=3pt, top=2pt, bottom=2pt]
resource "google\_storage\_bucket" "media" \{\\
\ \ name\ \ \ \ \ = "\$\{var.project\_id\}-\$\{var.app\_name\}-media"\\
\ \ project\ \ = var.project\_id\\[0.4em]
\ \ uniform\_bucket\_level\_access = true\\
\ \ public\_access\_prevention\ \ \ = "enforced"\\
\}
\end{extraitconf}

\section{Détection, livraison et fédération d'identité}

\textbf{Fenêtre et cadence de la règle R2.} L'extrait de
\texttt{terraform/\allowbreak modules/\allowbreak detection/\allowbreak main.tf}, lignes
103--114, montre la fenêtre SQL utilisée pour
compter les refus du filtrage applicatif.

\begin{extraitconf}[boxsep=1pt, left=3pt, right=3pt, top=2pt, bottom=2pt]
SELECT\\
\ \ JSON\_VALUE(json\_payload, "\$.httpRequest.remoteIp") AS src\_ip,\\
\ \ resource\_name AS service, COUNT(*) AS waf\_blocks\\
FROM \textasciigrave\$\{local.project\}.\$\{local.dataset\}.raw\_logs\textasciigrave\\
WHERE timestamp >= TIMESTAMP\_SUB(CURRENT\_TIMESTAMP(),\\
\ \ \ \ \ \ \ \ \ \ \ \ \ \ \ \ \ \ INTERVAL 15 MINUTE)\\
\ \ AND log\_source = "armor"\\
\ \ AND UPPER(JSON\_VALUE(json\_payload,\\
\ \ \ \ \ \ \ \ \ \ "\$.enforcedSecurityPolicy.outcome")) = "DENY"\\
GROUP BY src\_ip, service\\
WHERE f.waf\_blocks > 10 AND f.src\_ip IS NOT NULL
\end{extraitconf}

La ressource de requête planifiée fixe séparément \texttt{schedule = "every 5 minutes"} à la
ligne 410 : la fenêtre de quinze minutes n'est donc pas la cadence d'exécution.

\textbf{Déclenchement et permissions de la chaîne de livraison.} Le déclenchement filtré par
chemins et la déclaration des permissions figurent en tête du descripteur de la chaîne, et portent
deux décisions commentées au \cref{ch:realisation}.

\begin{extraitconf}[boxsep=1pt, left=3pt, right=3pt, top=2pt, bottom=2pt]
on:\\
\ \ push:\\
\ \ \ \ branches: [main]\\
\ \ \ \ paths: ["api/**", "dashboard/**", "frontend/**",\\
\ \ \ \ \ \ \ \ \ \ \ \ "scripts/**", "tests/**",\\
\ \ \ \ \ \ \ \ \ \ \ \ ".github/workflows/ci.yml"]\\
\ \ pull\_request:\\
\ \ \ \ branches: [main]\\[0.4em]
permissions:\\
\ \ contents: read\ \ \ \ \ \ \# lecture du depot uniquement\\
\ \ id-token: write\ \ \ \ \ \# federation d'identite : aucune cle
\end{extraitconf}

\textbf{Condition d'attribut de la fédération d'identité.} L'extrait suivant provient du module
\texttt{iam}. La condition porte \emph{à la fois} sur le dépôt et sur la référence de branche :
sans le second membre, toute branche du dépôt obtenait l'identité de déploiement
(\cref{subsec:wif}).

\begin{extraitconf}[boxsep=1pt, left=3pt, right=3pt, top=2pt, bottom=2pt]
attribute\_mapping = \{\\
\ \ "google.subject"\ \ \ \ \ \ \ = "assertion.sub"\\
\ \ "attribute.repository" = "assertion.repository"\\
\ \ "attribute.ref"\ \ \ \ \ \ \ \ = "assertion.ref"\\
\}\\[0.4em]
attribute\_condition = <<-EOT\\
\ \ assertion.repository == "\$\{var.github\_owner\}/\$\{var.repo\}"\\
\ \ \&\& assertion.ref == "\$\{var.allowed\_deploy\_ref\}"\\
EOT
\end{extraitconf}

\textbf{Discipline des exclusions d'analyse statique.} L'analyse statique est invoquée par appel
direct à l'outil, avec une version épinglée et une propagation fidèle du code de sortie. Les sept
exclusions sont nominatives et chacune porte, dans le descripteur, une justification ancrée sur un
fichier et une ligne. L'élision en tête d'identifiant remplace le préfixe d'espace de noms :
\texttt{...path-join-resolve-traversal} vaut pour
\texttt{javascript.\allowbreak lang.\allowbreak security.\allowbreak audit.\allowbreak
path-traversal.\allowbreak path-join-resolve-traversal}.

\begin{extraitconf}[boxsep=1pt, left=3pt, right=3pt, top=2pt, bottom=2pt]
pip install semgrep==1.173.0\\
semgrep scan --config=p/default --error \textbackslash\\
\ \ --exclude-rule=...express-path-join-resolve-traversal \textbackslash\\
\ \ --exclude-rule=...path-join-resolve-traversal \textbackslash\\
\ \ --exclude-rule=...detect-non-literal-regexp \textbackslash\\
\ \ --exclude-rule=...unsafe-formatstring \textbackslash\\
\ \ --exclude-rule=...dynamic-urllib-use-detected \textbackslash\\
\ \ --exclude-rule=...gcp-sql-database-ssl-insecure-value \textbackslash\\
\ \ --exclude-rule=...github-actions-mutable-action-tag\\[0.4em]
\# La regle bypass-tls-verification n'est PAS exclue :\\
\# chaque occurrence assumee porte une annotation nominative\\
\# nosemgrep sur sa propre ligne. Toute NOUVELLE occurrence\\
\# ailleurs fait echouer la chaine.
\end{extraitconf}

\section{Provisionnement, isolation et tableau de bord}

\textbf{Contrôles de la chaîne de provisionnement.} Le descripteur de la chaîne d'infrastructure
déclare la seule permission de lecture du dépôt : il ne peut structurellement pas s'authentifier
auprès du fournisseur cloud, et la validation s'exécute donc sans accès distant.

\begin{extraitconf}[boxsep=1pt, left=3pt, right=3pt, top=2pt, bottom=2pt]
permissions:\\
\ \ contents: read\ \ \ \ \ \ \# pas de id-token : aucun acces GCP\\[0.4em]
- run: terraform fmt -check -recursive\ \ \ \ \# bloquant\\
- run: for env in dev staging; do cd terraform/environments/\$env;\\
\ \ \ \ \ \ terraform init -backend=false -input=false; terraform validate; done\\[0.4em]
- uses: aquasecurity/trivy-action@v0.36.0\\
\ \ with:\\
\ \ \ \ scan-type: config\\
\ \ \ \ scan-ref:\ \ terraform/\\
\ \ \ \ severity:\ \ CRITICAL\\
\ \ \ \ exit-code: 1\ \ \ \ \ \ \ \ \ \ \ \ \ \ \# porte bloquante
\end{extraitconf}

\textbf{Isolation logique entre locataires.} Les comptes de base créés par le fournisseur héritent
d'un rôle privilégié : sans les ordres suivants, chaque application lit la base de l'autre.
L'extrait est partiel --- huit ordres, dont les six révocations qui portent le cloisonnement, quand
le corps en dénombre dix pour l'accueil complet du second locataire.

\begin{extraitconf}[boxsep=1pt, left=3pt, right=3pt, top=2pt, bottom=2pt]
REVOKE CONNECT ON DATABASE menal\_db FROM PUBLIC;\\
REVOKE CONNECT ON DATABASE menal\_db FROM app2\_user;\\
REVOKE CONNECT ON DATABASE app2\_db\ \ FROM PUBLIC;\\
REVOKE CONNECT ON DATABASE app2\_db\ \ FROM api\_user;\\
GRANT\ \ CONNECT ON DATABASE app2\_db\ \ TO\ \ \ app2\_user;\\
GRANT\ \ ALL ON SCHEMA public TO app2\_user;\\
REVOKE cloudsqlsuperuser FROM app2\_user;\ \ -- sinon l'acces\\
\ \ \ \ \ \ \ \ \ \ \ \ \ \ \ \ \ \ \ \ \ \ \ \ \ \ \ \ \ \ \ \ \ \ \ \ \ \ -- croise persiste\\
REVOKE cloudsqlsuperuser FROM api\_user;\ \ \ -- et de meme\\
\ \ \ \ \ \ \ \ \ \ \ \ \ \ \ \ \ \ \ \ \ \ \ \ \ \ \ \ \ \ \ \ \ \ \ \ \ \ -- dans l'autre sens
\end{extraitconf}

Les deux dernières révocations sont celles qui comptent : sans elles, l'accès croisé subsiste par
héritage, et leur symétrie n'est pas décorative. Omettre la seconde produirait une isolation
\textbf{à sens unique} --- la plateforme lirait la base du locataire sans la réciproque --- et c'est
la direction qui porte le risque : la base du locataire contient un identifiant national, celle de
la plateforme seulement des empreintes de mots de passe et un secret de second facteur chiffré.

\textbf{Chaîne de données d'une vue du tableau de bord.} Le motif suivant est identique dans les
dix éléments rendus au serveur : lecture du jeton dans le témoin de connexion, appel de l'API,
aucun accès direct à une source de données.

\begin{extraitconf}[boxsep=1pt, left=3pt, right=3pt, top=2pt, bottom=2pt]
// app/detections/page.tsx --- composant rendu au serveur\\
const token = cookies().get("token")?.value;\\
const data\ = await getDetections(token, 24, 200, tenant);\\[0.4em]
// lib/api.ts --- unique point de sortie HTTP\\
fetch(\textasciigrave\$\{apiUrl()\}\$\{path\}\textasciigrave, \{\\
\ \ headers: \{ Authorization: \textasciigrave Bearer \$\{token\}\textasciigrave{} \},\\
\ \ cache: "no-store",\\
\});
\end{extraitconf}

Aucun client d'entrepôt de données ni de base relationnelle ne figure parmi les dépendances de
production du tableau de bord : la propriété énoncée au \cref{ch:realisation} est donc
vérifiable par simple inventaire.

\section{Tableaux de référence de la réalisation}
\label{ann:tableaux-realisation}

Ces quatre tableaux sont des tableaux de \emph{consultation} : le \cref{ch:realisation} les
appelle sans les lire ligne à ligne.

{\scriptsize
\setlength{\tabcolsep}{3pt}%
\begin{longtable}{|P{3.05cm}|P{5.95cm}|P{5.80cm}|}
\caption[Familles de contrôles de la chaîne de livraison applicative]{Les huit familles de
contrôles de la chaîne de livraison applicative, dépôt vérifié le \dateref.}
\label{tab:familles-controles}\\
\hline
\textbf{Famille} & \textbf{État, outil et caractère bloquant} & \textbf{Destination du résultat} \\
\hline
\endfirsthead
\caption[]{Les huit familles de contrôles de la chaîne de livraison applicative (suite)} \\
\hline
\textbf{Famille} & \textbf{État, outil et caractère bloquant} & \textbf{Destination du résultat} \\
\hline
\endhead
Recherche de secrets versionnés & Réalisée --- détecteur de secrets, jeu de règles par défaut,
\textbf{bloquant} & Journaux d'exécution de la chaîne \\ \hline
Analyse statique du code & Réalisée --- analyseur statique, version épinglée, sept exclusions
nominatives, \textbf{bloquant} & Journaux d'exécution de la chaîne \\ \hline
Analyse de la description d'infrastructure & Réalisée --- analyseur de configuration,
\textbf{bloquant} sur gravité critique & Journaux d'exécution de la chaîne \\ \hline
Analyse de vulnérabilité des images & Réalisée --- analyseur d'image, quatre images,
\textbf{bloquant} sur gravité critique corrigeable & Journaux \textbf{et entrepôt de supervision}
(table des vulnérabilités) --- seule famille dont le résultat quitte la chaîne \\ \hline
Tests automatisés & Réalisée --- cadre Python (31 tests) et cadre JavaScript, \textbf{bloquant}
sur échec, non bloquant sur couverture & Rapport de test conservé 7~jours \\ \hline
Analyse de composition logicielle (dépendances) & \textbf{Arbitrée} & Risque capté en aval par
l'analyse d'image, qui inventorie les paquets de l'artefact livré ; ce chemin ne couvre ni les
dépendances de développement, ni les avis sans correctif amont, écartés par configuration.
Réévaluation : la première dépendance directe hors image \\ \hline
Test dynamique du service déployé & \textbf{Arbitrée} & La vérification post-déploiement établit
la disponibilité par le domaine public, jamais l'absence de vulnérabilité ; la validation dynamique
du filtrage applicatif est portée par les protocoles du \cref{ch:validation}. Réévaluation : une
homologation par un tiers \\ \hline
Signature d'artefact, provenance et nomenclature logicielle & \textbf{Arbitrée} (décision D07) &
Traçabilité atteinte par un chemin moins coûteux --- étiquette liée à l'empreinte du commit,
registre unique --- qui ne prouve pas que l'artefact déployé vient de la chaîne attendue.
Réévaluation : l'ouverture du socle à un éditeur tiers \\
\hline
\end{longtable}
}

\begin{table}[htbp]
\caption[Modèle de données de l'entrepôt de supervision]{Modèle de données de l'entrepôt de
supervision : dix tables, un écrivain chacune}
\label{tab:modele-donnees}
\centering
\scriptsize
\setlength{\tabcolsep}{3pt}
\renewcommand{\arraystretch}{1.15}
\begin{tabular}{|P{2.95cm}|P{3.55cm}|P{2.05cm}|P{3.75cm}|P{2.56cm}|}
\hline
\textbf{Table} & \textbf{Contenu} & \textbf{Partition et rétention} & \textbf{Écrivain autorisé} & \textbf{Clé d'application} \\
\hline
Journaux d'accès\newline \gcprod{access\_logs} & Une ligne par requête HTTP applicative & Jour · 90~j & Normalisation F4 (identité de déploiement) & \gcprod{service} \\ \hline
Journaux bruts\newline \gcprod{raw\_logs} & Charge JSON des quatre sources collectées & Jour · 90~j & Normalisation F4 & \gcprod{resource\_name} \\ \hline
Détections\newline \gcprod{detections} & Alertes R1--R7, tactique et technique ATT\&CK & Jour · \textbf{sans expiration} & Requêtes planifiées R1--R7 & \gcprod{service} \\ \hline
Événements de sécurité\newline \gcprod{security\_events} & Blocages du filtrage applicatif, déjà mitigés & Jour · \textbf{sans expiration} & Normalisation F4 & \textbf{absente} \\ \hline
Enrichissement\newline \gcprod{alert\_enrichment} & Technique rattachée, similarité, candidats de rang 2--3 & Jour · \textbf{sans expiration} & Tâche F5 --- droit limité à cette seule table & \textbf{absente} \\ \hline
Métriques\newline \gcprod{api\_metrics} & Agrégat horaire : volume, erreurs, latence, blocages & Jour · \textbf{sans expiration} & Normalisation F4 & \textbf{absente} \\ \hline
Verdicts\newline \gcprod{analyst\_verdicts} & Jugement humain, ajout seul, jamais de mise à jour & Jour · \textbf{sans expiration} & API de supervision --- cette seule table, rôle administrateur & \textbf{absente} \\ \hline
Vulnérabilités\newline \gcprod{cve\_findings} & CVE du dernier scan, catalogue d'exploitation, score EPSS & Jour · \textbf{sans expiration} & Chaîne d'intégration --- cette seule table, réécriture totale & sans objet \\ \hline
Vecteurs de référence\newline \gcprod{attack\_embeddings} & 872 vecteurs de 768 dimensions, catalogue ATT\&CK & Aucune & Chargement hors ligne, identité d'opérateur & sans objet \\ \hline
File d'attente ML\newline \gcprod{pending\_embeddings} & \textbf{Déclarée, jamais écrite ni lue} : vestige de F5 ; retrait décidé le \dateref, non appliqué & Aucune & Aucun & sans objet \\
\hline
\end{tabular}
\end{table}

{\scriptsize
\setlength{\tabcolsep}{3pt}%
\begin{longtable}{|P{2.60cm}|P{10.55cm}|P{1.65cm}|}
\caption[Cartographie des vues du tableau de bord]{Cartographie des vues du tableau de bord :
douze éléments rendus, dont dix au serveur, dépôt vérifié le \dateref.}
\label{tab:vues-dashboard}\\
\hline
\textbf{Vue} & \textbf{Apport à l'analyste et point d'entrée appelé} & \textbf{Filtre locataire} \\
\hline
\endfirsthead
\caption[]{Cartographie des vues du tableau de bord (suite)} \\
\hline
\textbf{Vue} & \textbf{Apport à l'analyste et point d'entrée appelé} & \textbf{Filtre locataire} \\
\hline
\endhead
Vue d'ensemble & Cinq indicateurs, répartition par gravité, courbe des requêtes, dix dernières
requêtes --- journaux, alertes et synthèse en parallèle & Partiel --- synthèse seule \\ \hline
Journal d'audit & Cent dernières requêtes reçues par l'API --- journaux & Non \\ \hline
Alertes & Requêtes en erreur, gravité dérivée du code de retour --- alertes, filtrées & Non \\ \hline
Détections & Deux cents détections des règles R1--R7 sur vingt-quatre heures --- détections & \textbf{Oui} \\ \hline
Incidents & Entités agrégées, score de tri, verdict posé --- incidents & \textbf{Oui} \\ \hline
Fiche d'incident & Jauge de score, tactiques rattachées, chronologie --- incident par entité,
200~lignes & \textbf{Oui} \\ \hline
Couverture ATT\&CK & Matrice de couverture par tactique --- couverture & \textbf{Oui} \\ \hline
Vulnérabilités & Priorisées : exploitation connue, observation, score de probabilité ---
vulnérabilités & Non \\ \hline
Santé des règles & Rendement et taux de faux positifs des sept règles --- santé des règles & Non \\ \hline
Enveloppe de mise en page & Navigation, bandeau, bascule de thème, indicateur de fraîcheur ---
aucun appel & Sans objet \\ \hline
Connexion \textit{(cliente)} & Formulaire en deux temps, mot de passe puis code temporel ---
relais de connexion puis de vérification & Sans objet \\ \hline
Sécurité du compte \textit{(cliente)} & Enrôlement et retrait du second facteur --- relais du
second facteur & Sans objet \\
\hline
\end{longtable}
}

\begin{table}[htbp]
\caption[Caractérisation quantifiée des deux logiciels écrits]{Caractérisation quantifiée des deux
logiciels écrits, dépôt vérifié le \dateref.}
\label{tab:logiciel-ecrit}
\centering
\scriptsize
\setlength{\tabcolsep}{3.5pt}
\renewcommand{\arraystretch}{1.10}
\begin{tabular}{|P{3.05cm}|P{6.10cm}|P{6.05cm}|}
\hline
 & \textbf{API de supervision} & \textbf{Tableau de bord} \\
\hline
Langage et cadre & Python 3.12, cadre applicatif asynchrone à validation typée &
TypeScript, cadre à rendu serveur \\ \hline
Volume écrit & 1\,855~lignes, 20~modules & 3\,408~lignes, 40~fichiers \\ \hline
Points d'entrée & 21, dont 9 pour la supervision & douze vues, dont dix rendues au serveur et
deux clientes, et sept routes serveur \\ \hline
Modèle de données & 4 tables, 3 migrations, 22 schémas de validation & 16 types partagés avec
l'API \\ \hline
Tests automatisés & 31 tests, exécutés à chaque livraison & 22 tests unitaires exécutés ;
22 tests de parcours écrits mais jamais exécutés par la chaîne \\ \hline
Couverture mesurée & non mesurée par la chaîne : l'outil est installé, l'option de mesure ne
l'est pas & \textbf{3,00~\% des lignes} --- 11~lignes couvertes sur 367, 22 fichiers
instrumentés --- pour un seuil de 70~\% déclaré et non appliqué \\ \hline
Dette assumée & aucune version portée par les points d'entrée : un client ne peut pas se lier
à un contrat stable & 375~lignes d'une interface antérieure, non déployée, conservées au
dépôt \\
\hline
\end{tabular}
\end{table}
